% Source-aligned English for printed p.156.
\SourcePage{156}

\LayerTarget{tohoku:proof:3.1.1:p2}
\noindent to an $\mathbf O$-module $F$ identify with the elements of $F(U)$.
Taking Theorem 1.10.1 into account gives the following.

\LayerTarget{tohoku:cor:3.1.1}
\indent\textsc{Corollary.} \textit{Every $\mathbf O$-module is isomorphic to a
sub-$\mathbf O$-module of an injective $\mathbf O$-module.}

\LayerTarget{tohoku:construction:godement-injective-product}
We give a direct proof of this corollary, due to Godement. For every $x\in X$,
let $M_x$ be an $\mathbf O(x)$-module, and define a sheaf $M$ on $X$ by
$M(U)=\displaystyle\prod_{x\in U}M_x$,
with the evident restriction maps and $\mathbf O(U)$-actions on $M(U)$. By
construction, $M$ is an $\mathbf O$-module on $X$ isomorphic to the product
of the $\mathbf O$-modules $M^x$ ($x\in X$), where $M^x(U)=M_x$ if $x\in U$
and $M^x(U)=0$ otherwise. It follows at once that for every $\mathbf O$-module
$F$, the homomorphisms from $F$ to $M$ identify with the families
$(u_x)_{x\in X}$ in which $u_x$ is an $\mathbf O(x)$-homomorphism from $F(x)$
to $M_x$. Hence:

\LayerTarget{tohoku:prop:3.1.2}
\indent\textsc{Proposition 3.1.2.} \textit{If, for every $x\in X$, $M_x$ is an
injective $\mathbf O_x$-module, then the product sheaf $M$ defined above is an
injective $\mathbf O$-module.}

Now let $F$ be any $\mathbf O$-module. It is classical, and also follows from
Theorem 1.10.1, that for every $x\in X$ the module $F(x)$ embeds into an
injective $\mathbf O_x$-module $M_x$. Consequently $F$ embeds into the
injective $\mathbf O$-module $M$ defined by the $M_x$.

\LayerTarget{tohoku:prop:3.1.3}
We also record the following for later use.

\indent\textsc{Proposition 3.1.3.} \textit{Let $M$ be an injective
$\mathbf O$-module on $X$, let $U$ be an open subset of $X$, and let
$\mathbf O_U$ and $M_U$ be the restrictions of $\mathbf O$ and $M$ to $U$.
Then $M_U$ is an injective $\mathbf O_U$-module.}

Clearly $M_U$ is an $\mathbf O_U$-module. Let $F$ be an $\mathbf O_U$-module,
let $G$ be a submodule of $F$, and let $u:G\to M_U$ be a
homomorphism. We show that $u$ extends to a homomorphism from $F$ to $M_U$.
For an $\mathbf O_U$-module $H$ on $U$, let $\overline H$ be the
$\mathbf O$-module obtained, in terms of étalé spaces, by ``extending $H$ by
zero on $\complement_U$''; see [4, XVII, prop. 1]. Giving an $\mathbf O_U$-homomorphism
$u:G\to M_U$ is then equivalent to giving an
$\mathbf O$-homomorphism $\overline G\to M$. Since
$\overline G$ is a submodule of $\overline F$ and $M$ is injective,
$\overline u$ extends to a homomorphism from $\overline F$ to $M$, which
induces the required homomorphism from $F$ to $M_U$. Note that Proposition
3.1.3 becomes \textit{false} if $U$ is assumed \textit{closed} rather than open. The same
argument proves:

\LayerTarget{tohoku:prop:3.1.4}
\indent\textsc{Proposition 3.1.4.} \textit{Let $M$ be an injective
$\mathbf O$-module on a closed subset $Y$ of $X$. Then the $\mathbf O$-module
$M^X$ which agrees with $M$ on $Y$ and is zero on $\complement Y$ is injective.}

\LayerTarget{tohoku:section:3.2}
\indent{\boldmath\textbf{3.2. Definition of $H^p_\Phi(X,F)$.}} Let $X$ be a topological space.
Write $\mathbf C^X$ for the abelian category $\mathbf C^{\mathbf Z}_X$ of abelian sheaves on
$X$. If $F$ is such a

\TohokuContinuation{157}

